\chapter{数据说明与实证设计}
\section{数据来源与当前可证实范围}
\subsection{公司数据的研究用途}
本文计划使用公司维护的中国股票分钟行情与配套日频、证券状态数据。分钟价格用于构造日度已实现方差，日收益用于估计条件风险关系，两者在相同证券与交易日历上对齐。主输入优先采用一分钟价格序列统一构造五分钟测度，存量五分钟数据仅在独立认证后作为对照，不以不同来源的历史数据静默拼补。

本版已读取原始数据位置索引及五份数据字典并记录文件哈希；本轮登记对象的实际字段结构均已取得，并已通过授权只读查询获取有界真实小样，公司行动严格版本视图也已完成有界读取。小样获取成功不等于验收通过：初步核对已识别收益来源口径差异及历史时点可见性证据不足，分钟会话规则仍在核验，逐股逐日覆盖尚未确认。字段存在及有限记录可读不能证明完整历史可用。表\ref{tab:data_status}说明后续验收事项；正式协议尚未锁定，具体日期、证券数量和覆盖统计仍待确认后填入。

\begin{table}[htbp]\centering\small
\caption{数据来源类别与验收状态}\label{tab:data_status}
\begin{tabularx}{\textwidth}{p{0.2\textwidth}YY}
\toprule 数据类别 & 已报告的用途与性质 & 实证前待确认事项\\\midrule
股票分钟价格 & L2成交价格快照派生，OHLC与区间量额 & 真实覆盖、区间标签、无成交规则、端点价格\\
股票日线 & 日收益和量额核对，含交易所参考前收 & 时间编码、参考价与公司行动一致性\\
公司行动 & 含市场可见时间的版本化视图，已取有界小样 & 历史覆盖、事件与条款版本联结、可见性证据\\
交易日历与状态 & 日历、停复牌、限价及质量记录 & 有效水位、缺口与哨兵值含义\\
历史成分与行业 & 带来源的成分、行业有效区间 & 发布证据、分类版本、历史修订\\
证券身份 & 代码、市场、上市与退市字段 & 当前快照限制与跨代码主体歧义\\\bottomrule
\end{tabularx}
\end{table}

本研究不需要客户账户、公司交易策略、持仓或逐笔委托数据。内部连接细节由独立的数据适配文档管理，论文仅说明可披露的数据类别、频率、时段与处理方法。最终是否公开派生RV和风险边，也须获得明确的数据授权。

\subsection{三个不同层次的时点保证}
本文将时点问题分为三个层次。第一，算法时间顺序，要求起点$t$的特征、图和预测不能读取$t$之后的观测。第二，事件及分类可见性，要求使用可证明在当时公开的版本，而不是将后来的修订提前用于历史。第三，行情本身的历史版本与实际入库时刻，即当前历史库是否还原了当时实际可读的记录。

本稿计划对第一层实施严格测试，并在公司行动等具备证据的部分使用第二层。第三层目前尚未取得完整证明。因此，即使所有算法都按时间顺序运行，也只能在相应数据快照上声称时间外预测，不能自动声称完整复原了历史线上可得信息。程序运行时间不是历史发布时点，后补的历史记录也不能仅凭业务日期获得当时可见的资格。

原始字典声明，交易日$D$的分钟数据最早在$D+1$日北京时间05:00之后、日线最早在06:10之后，且相应任务成功后才可用。这是待实测确认的文档条件，不是已验证的服务时效保证，也不能据此称完整输入在$D$日收盘时即可读取。第三章中的$t$首先表示最后纳入观测的交易日；真正发出预测的时刻还须晚于全部输入的可用时刻。正式协议须确认该时刻是否早于目标首日的交易会话；若不能满足，需在测试前调整预测起点与标签对齐，而不能把延迟到达的数据提前使用。

\section{样本选择与历史覆盖}
\subsection{样本形成规则}
开题提出约30只具有流动性和行业代表性的沪深300股票作为主样本。本文保留这一规模目标，但将形成日设在首个预测模型训练起点之前：先用历史预热段确定可估计风险关系的候选证券，在形成日根据当时可证明的成分资格、此前成交额、历史可用性和行业信息选样，然后固定集合开展后续训练、验证和测试。

这种安排不使用测试期是否持续交易、是否退市或模型预测效果进行反向选择。样本形成后发生的停牌、成分退出和证券身份变化均作为后续数据状态记录，不因“不方便”而将该股票从整个历史中删除。形成日前的预热数据用于选择和初始化，不作为已经对该固定样本开展在线预测的证据。

如果历史成分的公布时间不可证实，则研究应降级为明确限定的固定股票样本，不再声称严格的历史沪深300可投资成分筛选。若完整跨代码证券身份不可得，则采用交易所、代码与上市身份区间构造临时键，禁止仅凭股票名称串接不同主体。

\subsection{历史行业与共同信息}
行业匹配使用有效区间和分类版本共同确定。名称含某一版本年份的表并不代表全部历史都采用该版本，追溯性分类也不自动等于当年已知信息。不能以当前行业成员替代历史成员，并将由此形成的行业平均值当作完全无前视的控制。

在历史行业无法认证时，匹配组统一删去行业变量，但仍可使用剔除目标股票的样本平均RV：
\begin{equation}\label{eq:market_control}
M_{-i,t}=\frac{1}{|\mathcal V_t\setminus\{i\}|}
\sum_{j\in\mathcal V_t\setminus\{i\}}\RV_{j,t},
\end{equation}
其中$\mathcal V_t$为当日测度有效的股票集合，并同时记录有效数量。该变量是样本中个股RV的平均值，不是市场指数RV；后者还受到成分协方差及权重影响。市场指数分钟数据尚未认证，不阻塞采用明确定义的共同信息对照。

\subsection{覆盖统计与样本长度}
较长历史用于容纳风险网络初始化、预测模型训练、验证和测试。初始预算候选为750个交易日的风险预热、750个训练起点、250个验证起点和500个测试起点。它们是开发规划，不是已查得的覆盖或通用最低样本定理；实际还要扣除缺失、特征回看和标签成熟造成的损失。

覆盖统计至少包含原始分钟记录、有效五分钟端点、完整股票日、连续HAR窗口及一、五、二十二日成熟标签。仅报告最早与最晚日期，不足以证明中间没有缺口；同日30只股票也不能直接视作30个独立市场时点。若真实历史不足，须在最终测试之前调整规模或切分，并明确统计精度的限制。

\begin{table}[htbp]\centering\small
\caption{正式样本待验收信息}\label{tab:sample_pending}
\begin{tabularx}{\textwidth}{p{0.34\textwidth}Y}
\toprule 项目 & 本版状态\\\midrule
样本股票与证券身份列表 & \Pending{实际名单和形成日}\\
分钟原始区间与有效交易日 & \Pending{逐股逐日覆盖统计}\\
收益调整证据与公司行动覆盖 & \Pending{核对结果和版本限制}\\
训练、验证与测试区间 & \Pending{冻结日期与数据版本}\\
各期限有效标签及排除原因 & \Pending{样本损失分解}\\\bottomrule
\end{tabularx}
\end{table}

\section{数据预处理与测度验收}
\subsection{时间编码与会话边界}
已报告的源合同存在重要差异：一分钟导出记录的UTC样式后缀可能只是裸北京时间的格式标记，而旧五分钟源的同一后缀才表示真实UTC。因此，不能对所有表统一使用“转换到北京时间”的逻辑。

实际接入按来源与传输接口分别建立时间profile，比较原始字符串或epoch、数据库精度、会话时区和解码后的开盘、午休及收盘位置。一旦内部统一为带时区的北京时间，禁止再次解释导致八小时偏移。日线与指数数据也单独认证，不直接继承分钟表的时间规则。

五分钟价格路径在上午和下午分别建立。开盘集合竞价单独成根时，应明确其作为锚点还是进入区间汇总，并避免重复计算；下午首价的真实时点需记录。价格保持仅限于已证实的无成交区间，不能用后续价格反填缺失，也不能跨休市期间生成不存在的成交区间。

字典另声明一分钟记录采用右端标签、09:30开盘集合竞价单列、15:00记录包含收盘集合竞价。这些声明需要用授权小样及真实schema核对后才进入有效采样规则；本稿不将它们当作已经观察到的数据库事实。

\subsection{价格、量额与异常状态}
数据检查包括同键重复、冲突版本、价格正值、OHLC约束、量额单位和交易状态。完全相同的重复可以按固定规则去重；同键不同值的冲突需保留证据，不按文件最后一行直接覆盖。区间股数与区间成交额先与日线核对口径，再评价差值。

真实极端波动、涨跌停及高成交额本身不是数据错误。本文主标签不按全样本分位数缩尾，也不为降低MSE删除真实高波动日。停牌、无成交和数据缺口分开处理：停牌不自动等于零方差；没有记录不自动证明没有成交；零成交量下重复的最新成交价是否可靠，需由源生成逻辑和样本确认。

\subsection{收益调整与公司行动}
候选参考价收益式\eqref{eq:reference_return}需对普通日、除权除息日及其他特殊事件进行核验。若采用价格调整因子$A_{i,t}$，定义
\begin{equation}\label{eq:adjusted_return}
R^{\mathrm{adj}}_{i,t}=\log\frac{P_{i,t}A_{i,t}}{P_{i,t-1}A_{i,t-1}}.
\end{equation}
它与参考价收益一致所需的因子方向为
\begin{equation}\label{eq:factor_direction}
\frac{A_{i,t}}{A_{i,t-1}}=\frac{P_{i,t-1}}{B_{i,t}}.
\end{equation}
数据库字段名为ratio并不保证恰好定义为式\eqref{eq:factor_direction}，必须核对实际说明与样例。现金分红、股份变动和权利到账还会影响财富回报，因此价格调整恒等式不构成总回报认证。

版本化公司行动需同时满足两个时间条件：使用在预测截止之前已知的事件版本，并在经济生效时处理相应价格机械变化。不能在公告当天提前实施未来除权调整，也不能只获取最新事件快照回填过去。事件覆盖表用于区别“确实无事件”和“尚未覆盖”，没有查到事件行不能自动得出前者。

\subsection{数据质量与模型有效样本}
每个股票日保存测度有效性、特征有效性和标签有效性。特征缺失不应自动删除其他股票的当日记录；图边失效也不能成为主模型独占的删样理由。所有模型首先在同一标签域内生成预测，对失败使用预设兜底，并分别报告失败率和兜底占比。

在实现验收中，应选取真实完整股票日手算RV，并核对至少一种特殊会话或公司行动情形。纯构造数据用于测试公式和边界，但不代替真实数据清洗验证，更不进入主市场结果表。

\section{实证模型与贡献归因}
\subsection{模型体系}
\begin{table}[htbp]\centering\small
\caption{主比较模型及其作用}\label{tab:models}
\begin{tabularx}{\textwidth}{p{0.23\textwidth}YY}
\toprule 模型组 & 主要信息/结构 & 对应比较目的\\\midrule
传统HAR、HARQ & 自身多尺度历史、四次变差 & 强计量基准\\
匹配正值HAR、自身MLP & 相同缩放、链接与优化 & 损失和非线性容量\\
公共信息增强模型 & 市场/可用行业/边际风险 & 排除共同风险解释\\
风险汇总MLP & 相同公共信息和风险汇总 & 核心关系增量基线\\
Ridge-VHAR & 正则化横截面HAR输入 & 非图跨股票竞争者\\
GHAR、GLASSO-GNN & 线性/非线性统计图 & 图骨干与一般关联\\
DCC相关边模型 & 相同动态相关来源 & 排除仅有动态性收益\\
CoVaR二值/标量模型 & 固定支持、不同边输入 & 拓扑与边向量区分\\
底层参数边模型 & 当前及滞后原始参数 & 风险表征归纳偏置\\
完整边特征模型 & 相同自身/汇总+逐边消息 & 主确认性模型\\\bottomrule
\end{tabularx}
\end{table}
传统基线用于文献比较，统一架构组用于内部归因。GLASSO原文复现与统一骨干变体单独注册；忠实复现不能因使用相近名字就省略输入区间、归一化和损失的核对。相关构图使用何种收益、是否保留负关联，以及如何匹配边预算，均写入配置。

\subsection{主实验与消融顺序}
实证依次开展四个层次。第一层比较自身模型在MSE与QLIKE训练下的表现。第二层依次加入公共信息、风险汇总和逐边分支，检验完整模型是否超过匹配风险汇总基线。第三层保持节点信息、连接预算和骨干一致，替换动态相关、CoVaR拓扑、标量边与向量边。第四层通过尺度标准化、底层参数和关系破坏实验，检验改善是否依赖特定风险语义和邻居对应关系。

关系破坏包括固定拓扑重配入边特征、保度或近似保度重连、对称化及使用滞后图。作为正式消融时重新训练全部模型；仅在推理时改变边的结果则作为敏感性分析，不能与重训消融混合。随机重连改变了哪些结构、转置是否改变入度，均需说明。

动态性另拆为依赖状态更新与边际尺度更新：固定其中一部分、仅更新另一部分，再与全部更新比较。这样避免将“波动尺度在变”全部解释为“依赖网络在变”。上行通道若由数学性质决定重复，就不进行两个独立信号的贡献宣称。

\section{样本外训练与预测协议}
\subsection{时间切分与标签成熟}
设$c$为一次模型训练截止日，允许进入拟合的样本满足
\begin{equation}\label{eq:training_set}
\mathcal T_c^{(h)}=\{(i,s):\ \mathrm{end}(s,h)\leq c,
\ \mathrm{available}(y_{i,s}^{(h)})\leq c,\ \mathrm{valid}_{i,s,h}=1\}.
\end{equation}
这里的available只有在具备真实可得时间证据时才直接使用；否则明确采用固定历史快照下的成熟规则，并披露不能验证入库延迟。仅起点$s$已经过去，并不表示未来22日标签已经成熟。

训练、验证和最终测试按连续时间切分，不能随机打散股票日或图快照。跨边界目标窗口按最大期限22日的共同规则隔离，另可报告期限专属规则。每次拟合保存训练起点范围、目标终点范围和实际可用日期，使标签泄漏可以被自动检测。

\subsection{滚动更新与调参}
开发段为每个学习模型族提供至多12组候选设置，以两个固定随机种子选择结构与训练轮数。默认开发起点为一层图、隐藏维度32、学习率$10^{-3}$、权重衰减$10^{-4}$和dropout 0.1，均是待验证的工程参数，而非已观察到的最优值。

正式测试约每21个交易日重训，使用开发期确定的训练轮数，并按同一规则纳入已成熟的过去样本。不得根据最终测试表现修改图密度、删除年份、改变主要期限或重选胜出的随机种子。五个固定种子用于反映优化随机性，确定性基线不通过复制运行变成更多独立市场证据。

\subsection{历史不变性与失败规则}
将$t$之后的行情修改或追加后，在同一已冻结证券集合和随机状态下，截至$t$的特征、图和预测应保持不变。该测试同时覆盖价格处理、标准化、构图和训练数据索引。当前时点重估参数可以改善当前状态，但不能覆盖更早起点已归档的图。

图估计失败时优先使用年龄未超限的合法参数重新过滤；不可用则缺边。预测层保留无图回退。若无图模型也不可用，使用预设的近期有效RV或训练历史均值，并记录原因。主评价保留这类操作后预测，防止模型通过不输出困难样本改善结果；纯成功子集仅作附表。

\section{评价指标与统计检验}
\subsection{日期级损失与相对改善}
对模型$m$，设$I_t^{(h)}$为所有模型共同的目标有效股票集合，$S$为随机种子数，定义日期级模型族损失
\begin{equation}\label{eq:daily_loss}
\ell_{m,t}^{(h)}=\frac1S\sum_{s=1}^{S}\frac1{|I_t^{(h)}|}
\sum_{i\in I_t^{(h)}}L\big(y_{i,t}^{(h)},\widehat y_{m,s,i,t}^{(h)}\big).
\end{equation}
首先对同日截面与种子取平均，再以日期为统计单位。种子平均损失不等于先平均预测再计算损失；后者属于另一个集成模型，需要单独登记。

记基线为$b$、主模型为$m$，定义
\begin{equation}\label{eq:difference}
d_t=\ell_{b,t}-\ell_{m,t},\qquad
\Delta_{b,m}=\frac1T\sum_t d_t,\qquad
\mathrm{Gain}_{b,m}=1-\frac{\bar\ell_m}{\bar\ell_b}.
\end{equation}
正的$\Delta$或Gain表示主模型损失较低，但Gain是损失的相对降低，不是收益率或交易利润。主指标为QLIKE，MSE、RMSE和MAE辅助报告，所有表明确单位和有效样本数。

\subsection{区块置信区间与 DM 检验}
日期级差值保留了同日股票共同冲击。主置信区间采用成对日期区块重采样，同时保留不同模型在同一天的对应关系。候选设置为2000次重采样、平均块长50个交易日，并以25和100作敏感性；这些设置在测试之前冻结。

DM比较以损失差均值为对象\cite{diebold1995}。采用HAC长期方差估计$\widehat\Omega$，统计量为
\begin{equation}\label{eq:dm}
\mathrm{DM}=\frac{\bar d}{\sqrt{\widehat\Omega/T}},\qquad
\widehat\Omega=\widehat\gamma_0+2\sum_{k=1}^{K}\left(1-\frac{k}{K+1}\right)\widehat\gamma_k.
\end{equation}
候选带宽$K=\max(h-1,20)$并作敏感性分析。多日目标重叠、滚动参数估计和厚尾损失差会影响有限样本推断，因此报告效应量、置信区间及适用限制，不只报告一个$p$值。

\subsection{多重比较、MCS 与市场状态}
主比较预设为完整模型对匹配风险汇总、公共信息、动态相关、标量CoVaR和底层参数模型等少量差值，并对比较族实施Holm校正。MCS在相同日期损失矩阵上计算90\%与95\%集合\cite{hansen2011}。进入集合并不表示显著优于集合内其他模型。

高风险状态由起点$t$已经可得的市场汇总RV与训练确定的阈值比较得到，例如开发历史的90\%分位数。未来实际极端日可以用于事后误差剖析，但不进入实时状态变量。高、低状态差异应直接检验差值之差，不能以“高状态显著、低状态不显著”替代组间差异检验。

\section{稳健性、计算成本与本章小结}
预设稳健性覆盖风险估计窗口、图稀疏度、采样频率、点条件与事件条件、边际模型及样本规模。所有参数不做无目的全笛卡尔积，而围绕核心结果分层执行。图估计对所有模型和种子共享缓存；报告边际与copula估计时间、模型训练时间、推理时间、失败率和内存需求，以评价复杂度与预测收益之间的取舍。

本章将数据位置转化为可验收的数据协议，并将研究问题转化为有明确对照、时间范围和统计单位的实验。下一步先验证真实分钟形态与收益口径，再锁定样本和切分。只有完成这些环节，第五章中的实证数字才具有明确的数据与方法含义。
